% !TeX root = ../main.tex
\section{Simply typed lambda calculus}
\label{sec:stlc}

The untyped calculus accepts terms such as $\mathit{TRUE}\ 3$ or $x\, x$. The
\textbf{simply typed lambda calculus} (STLC) adds the smallest type discipline:
every term has a type, and an application is legal only when the types
match. Its rules are the rules of LETREC (Section~\ref{sec:letrec}) for
variables, functions, application and \texttt{if}; LETREC adds numbers,
\texttt{let} and recursion.

\subsection{Types and terms}

Types are built from two base types and the function type:
\[
\tau ::= \tyint \mid \tybool \mid \tau \to \tau.
\]
Terms are those of the untyped calculus, with an annotation on each parameter,
plus literals and conditionals:
\[
e ::= x \mid \lambda x{:}\tau.\, e \mid e\, e \mid n \mid \ift{e}{e}{e}.
\]

\subsection{The typing judgment}

\begin{mathpar}
\inferrule*[right=T-Var]{x : \tau \in \Gamma}{\Gamma \vdash x : \tau}
\and
\inferrule*[right=T-Abs]{\Gamma, x : \tau_1 \vdash e : \tau_2}{\Gamma \vdash \lambda x{:}\tau_1.\, e : \tau_1 \to \tau_2}
\and
\inferrule*[right=T-App]{\Gamma \vdash e_1 : \tau_1 \to \tau_2 \\ \Gamma \vdash e_2 : \tau_1}{\Gamma \vdash e_1\, e_2 : \tau_2}
\and
\inferrule*[right=T-If]{\Gamma \vdash e_1 : \tybool \\ \Gamma \vdash e_2 : \tau \\ \Gamma \vdash e_3 : \tau}{\Gamma \vdash \ift{e_1}{e_2}{e_3} : \tau}
\end{mathpar}
The rules are syntax-directed: for each form of term exactly one rule
applies, so checking a term never requires a guess. \textsc{T-App} rejects
$3\, (\lambda x{:}\tyint.\, x)$, because 3 does not have a function type, and
$(\lambda x{:}\tyint.\, x)\ \mathit{true}$, because $\tyint \ne \tybool$.
Figure~\ref{fig:stlc-deriv} shows a derivation.

\begin{figure}[h]
\centering
\begin{mathpar}
\inferrule*[right=T-Abs]
  {\inferrule*[right=T-Abs]
    {\inferrule*[right=T-Var]{x : \tyint \in \{x : \tyint, y : \tybool\}}{x : \tyint, y : \tybool \vdash x : \tyint}}
    {x : \tyint \vdash \lambda y{:}\tybool.\, x : \tybool \to \tyint}}
  {\vdash \lambda x{:}\tyint.\, \lambda y{:}\tybool.\, x : \tyint \to \tybool \to \tyint}
\end{mathpar}
\caption{Typing derivation of $\lambda x{:}\tyint.\, \lambda y{:}\tybool.\, x$.}
\label{fig:stlc-deriv}
\end{figure}

\subsection{Capture-avoiding substitution}

$\beta$-reduction substitutes the argument into the body, $e[x := v]$, and the
substitution must not capture variables. With $e = \lambda y.\, x\, y$ and
$v = y$, replacing $x$ by $y$ textually gives $\lambda y.\, y\, y$, where the $y$
that was free has been captured by $\lambda y$. The fix is to rename the bound
variable first ($\lambda z.\, x\, z$, which is $\alpha$-equivalent) and then
substitute: $\lambda z.\, y\, z$. Compilers often avoid the problem by replacing
names with de Bruijn indices, the distance to the binder.

\subsection{Strong normalization}

\begin{theorem}[Strong normalization]
If $\vdash e : \tau$, every reduction sequence that starts at $e$ ends in a
normal form.
\end{theorem}

So every well-typed STLC program terminates, and STLC is not Turing complete.
The reason is that self-application has no type. For $x\, x$ to be typed, $x$
needs a type $\alpha$ with $\alpha = \alpha \to \beta$, and no finite type satisfies
that equation, because the right side is always bigger than the left. The
combinator $Y$ is built from self-applications, so it has no type either.
Section~\ref{sec:inference} meets the same equation as the occurs check of
\texttt{proc (f) zero?((f f))}, and Section~\ref{sec:recursive} adds recursive
types to give it a solution. LETREC recovers recursion directly with
\texttt{letrec}, and that is why it can loop.

\subsection{Propositions as types}

Read as logical formulas, the types of STLC are propositions of
intuitionistic logic, and the well-typed terms are their proofs
(Table~\ref{tab:curry-howard}). \textsc{T-App} is modus ponens: from a proof of
$P \Rightarrow Q$ and a proof of $P$ it gives a proof of $Q$. \textsc{T-Abs} is
conditional proof: to prove $P \Rightarrow Q$, assume $P$ and derive $Q$.

\begin{table}[h]
\centering
\small
\begin{tabular}{l l}
  \toprule
  Logic & Types \\
  \midrule
  proposition & type \\
  proof of $P$ & term of type $P$ \\
  implication $P \Rightarrow Q$ & function type $P \to Q$ \\
  conjunction $P \wedge Q$ & product type $P \times Q$ \\
  modus ponens & application \\
  \bottomrule
\end{tabular}
\caption{The Curry--Howard correspondence.}
\label{tab:curry-howard}
\end{table}

The term $\lambda x{:}\tau_1.\, \lambda y{:}\tau_2.\, x$ of Figure~\ref{fig:stlc-deriv}
is a proof of $P \Rightarrow (Q \Rightarrow P)$. Proof assistants such as Coq and
Agda are built on this correspondence: checking the type of a program is
checking a proof.

\subsection{Exercises}

\begin{exercise}
Is $\lambda x{:}(\tyint \to \tyint).\, x\, x$ well typed?
\end{exercise}

\paragraph{Solution.} No. \textsc{T-App} on $x\, x$ needs the type of the operator,
$\tyint \to \tyint$, to have as argument type the type of the operand,
$\tyint \to \tyint$. But the argument type is $\tyint$, and $\tyint \ne \tyint \to \tyint$.

\begin{exercise}
Perform $(\lambda y.\, x\, (\lambda x.\, x\, y))[x := y\, z]$.
\end{exercise}

\paragraph{Solution.} The outer $\lambda y$ would capture the $y$ of $y\, z$, so rename it
to $w$: $\lambda w.\, x\, (\lambda x.\, x\, w)$. Inside $\lambda x$ the variable $x$ is
bound, so only the first $x$ is replaced: $\lambda w.\, (y\, z)\, (\lambda x.\, x\, w)$.

\begin{exercise}
Write a term of type $(\tau_1 \to \tau_2) \to (\tau_2 \to \tau_3) \to (\tau_1 \to \tau_3)$.
Which logical rule does it prove?
\end{exercise}

\paragraph{Solution.} $\lambda f{:}\tau_1 \to \tau_2.\, \lambda g{:}\tau_2 \to \tau_3.\, \lambda x{:}\tau_1.\, g\,(f\, x)$,
function composition. It proves that implication is transitive: if
$P \Rightarrow Q$ and $Q \Rightarrow R$, then $P \Rightarrow R$.
